\chapter{DESENVOLVIMENTO}
\label{cha:desenvolvimento}

Este capítulo apresenta o desenvolvimento da arquitetura proposta, destacando as decisões técnicas adotadas e os experimentos conduzidos.

\section{Arquitetura Proposta}

São exploradas duas abordagens: (i) uma baseada em extração de características e aprendizado de máquina tradicional, e (ii) outra baseada em aprendizado profundo com redes neurais convolucionais.

\begin{enumerate}
    \item A primeira utiliza descritores (HOG e LBP) seguidos de classificadores (SVM e RF).
    \item A segunda aplica CNNs diretamente nas imagens de rastreamento ocular.
\end{enumerate}

\begin{figure}[h!]
\centering
\caption{Fluxo comparativo entre as abordagens de AM tradicional e Aprendizado Profundo.}
\includegraphics[width=13cm]{Imagens/metodologia_am.png}
\label{fig:compara}
\fonte{Adaptado de \citeonline{img_trad_dl}.}
\end{figure}

\section{Base de Dados}

Foi utilizada a base DOVES, composta por imagens de estímulo e dados de rastreamento ocular de 29 voluntários saudáveis.

\begin{itemize}
    \item 101 imagens por participante (totalizando 2.929 amostras)
    \item Cada imagem tem dimensão de 1024x768 pixels, em escala de cinza
    \item Os movimentos foram coletados com precisão sub-milimétrica a 200Hz
\end{itemize}

\section{Extração de Características}

\subsection{HOG (Histograma de Gradientes Orientados)}

Utilizado para capturar bordas e texturas nos rastros oculares. Resultado: vetor de 76.320 dimensões.

\subsection{LBP (Padrão Binário Local)}

Captura microtexturas da imagem. Resultado: vetor de 130.626 dimensões.

\section{Modelos de Aprendizado de Máquina}

\subsection{SVM (Máquinas de Vetores de Suporte)}

Classificador binário estendido para multi-classe via abordagem Um-Contra-Um. Foram utilizados diferentes kernels.

\subsection{Random Forest (Floresta Aleatória)}

Classificador do tipo ensemble baseado em múltiplas árvores de decisão. Utilizado pelo seu bom desempenho em conjuntos de alta dimensionalidade.

\subsection{CNN (Redes Neurais Convolucionais)}

Aplicadas via transferência de aprendizado. Duas arquiteturas foram utilizadas:

\begin{itemize}
    \item VGG-19
    \item ResNet-50
\end{itemize}

\section{Métricas de Avaliação}

As métricas utilizadas foram:

\begin{itemize}
    \item FAR – False Acceptance Rate
    \item FRR – False Rejection Rate
    \item ACC – Acurácia média
\end{itemize}

Os cálculos foram feitos a partir das matrizes de confusão para o cenário multi-classe.

